\subsection*{Host Code}

The host side uses the same \texttt{make\_spandata} / \texttt{.view()} pattern from earlier chapters:

\begin{lstlisting}[language=croqtile,caption={Host Code}]
int main() {
  auto lhs = croq::make_spandata<croq::s32>(128, 256);
  auto rhs = croq::make_spandata<croq::s32>(256, 256);
  lhs.fill_random(-10, 10);
  rhs.fill_random(-10, 10);
  auto res = matmul(lhs.view(), rhs.view());
  for (size_t i = 0; i < res.shape()[0]; ++i)
    for (size_t j = 0; j < res.shape()[1]; ++j) {
      int ref = 0;
      for (size_t k = 0; k < lhs.shape()[1]; ++k)
        ref += lhs[i][k] * rhs[k][j];
      croq::croq_assert(ref == res[i][j], "values are not equal.");
    }
  std::cout << "Test Passed\n" << std::endl;
}
\end{lstlisting}

Nothing new here. The host program stays boring so the Croqtile function stays the star.
